% =====================================================
% Baby model: when does optimal advertising rise with
% a per-unit excise tax? Worked-out math + scaffolding.
% Prose is intentionally terse — fill in for the paper.
% =====================================================

\section*{When Does Optimal Advertising Rise with a Tax?
A Sign Condition}

% [SCAFFOLDING / NOTES FOR PAPER PROSE]
%
% This section answers: under what condition does a
% profit-maximizing firm RAISE advertising in response
% to a per-unit excise tax on its product?
%
% Frame against the companion paper (sb_incidence draft,
% Section "Theory"), which assumes per-customer demand
% D(p) depends only on price, not on advertising A.
% Under that assumption a tax reduces optimal ads. We
% show the result hinges on that specific assumption.
% Relaxing it — allowing ads and price to enter demand
% jointly — gives a clean sign condition on the cross-
% partial D_{pA} that flips the comparative static.

\paragraph{Setup.} A monopolist sells a good with
demand $D(p, A)$, where $p$ is the consumer price and
$A \geq 0$ is total advertising (in dollars). Marginal
cost is $c$, and the government levies a per-unit
specific tax $t \geq 0$. The firm's profit is
\begin{equation}\label{eq:profit}
    \pi(p, A; t)
    \;=\; (p - c - t)\, D(p, A) \;-\; A.
\end{equation}
We assume $D_p < 0$, $D_A > 0$, $D_{AA} < 0$ (diminishing
returns to advertising), and the regularity conditions
needed for an interior optimum.

\begin{remark}[Relation to companion model]
The companion paper (sb\_incidence,
\texttt{paper/draft.tex}, the "Model" section
introducing $\pi^{\text{pre-tax}} = p D(p) A - c(A)$)
writes profit as
$\pi = (p-t) \, D(p) \, A - c(A)$,
where $A$ is the number of customers acquired (one ad
unit per customer) and per-customer demand $D(p)$
depends only on price (its Assumption~3). This is the
multiplicatively separable case $D(p, A) = D(p) \cdot A$,
with $D_{pA} = D'(p) < 0$ and a particular cost function.
The present analysis nests that as a special case and
shows it sits strictly on the $dA^*/dt < 0$ side of the
sign condition.
\end{remark}

\paragraph{First-order conditions.} At an interior
optimum, $(\pi_p, \pi_A) = (0, 0)$:
\begin{align}
    \pi_p &: \quad D \;+\; (p - c - t)\, D_p \;=\; 0,
        \label{eq:foc_p} \\
    \pi_A &: \quad (p - c - t)\, D_A \;-\; 1 \;=\; 0.
        \label{eq:foc_A}
\end{align}
Define the per-unit margin $m \equiv p - c - t$. The
price FOC is the standard Lerner condition
$m = -D / D_p$; the advertising FOC is Dorfman--Steiner
$m \, D_A = 1$.

\paragraph{Comparative static.} Let $p^*(t)$ and $A^*(t)$
be the optimal choices. Define the pass-through
$\rho \equiv dp^*/dt \in (0, 1)$ (typically less than one
under standard demand curvature). Total-differentiating
the advertising FOC \eqref{eq:foc_A} with respect to $t$
at the optimum:
\begin{equation}\label{eq:dAdt_unsolved}
    (\rho - 1)\, D_A
    \;+\; m \,\bigl[\, D_{Ap}\, \rho \;+\; D_{AA}\, A_t \,\bigr]
    \;=\; 0,
\end{equation}
where $A_t \equiv dA^*/dt$. Solving:
\begin{equation}\label{eq:dAdt}
    \boxed{\quad
    \frac{dA^*}{dt}
    \;=\;
    \frac{(\rho - 1)\, D_A \;+\; m \, \rho \, D_{Ap}}
         {-\, m \, D_{AA}}.
    \quad}
\end{equation}
The denominator $-m D_{AA} > 0$ by the second-order
condition for $A$, so $\mathrm{sign}(A_t)$ is the sign
of the numerator.

\begin{remark}[Uniqueness of the optimum]
\label{rem:uniqueness}
Conditions $D_p < 0$ (Marshall's Second Law) and
$D_{AA} < 0$ (diminishing returns to ads) make
$\pi_{pp}, \pi_{AA} < 0$ but are not enough for joint
concavity in $(p, A)$. We also need the off-diagonal
Hessian term not to dominate:
$\pi_{pp}\, \pi_{AA} - \pi_{pA}^2 > 0$, where
$\pi_{pA} = D_A + m\, D_{pA}$. The same cross-partial
$D_{pA}$ that drives the sign condition in
Proposition~\ref{prop:sign} can, if large enough,
break joint concavity. We assume it does not — i.e.,
$D_{pA}$ is positive enough to satisfy
\eqref{eq:sign_condition} but bounded above by
$\sqrt{\pi_{pp}\, \pi_{AA}}/m - D_A/m$ to preserve a
unique interior optimum.
\end{remark}

\begin{proposition}[Sign condition]\label{prop:sign}
Optimal advertising rises with the per-unit tax,
$dA^*/dt > 0$, if and only if
\begin{equation}\label{eq:sign_condition}
    D_{Ap}
    \;>\;
    \frac{(1 - \rho)\, D_A}{m \, \rho}.
\end{equation}
The RHS is strictly positive whenever pass-through is
less than full ($\rho < 1$) and the margin is positive
($m > 0$). Hence the firm advertises more after a tax
\emph{only if} $D_{Ap}$ is positive and large enough.
\end{proposition}

% [SCAFFOLDING: economics of D_{pA} > 0]
% D_{pA} > 0 means the marginal demand response to
% advertising is larger when the price is higher. This
% is the persuasive view: ads matter more for retention
% / brand loyalty / habit when the consumer has a stronger
% reason to defect (price is high). Becker-Murphy (1993)
% and Galbraith are the canonical references.
%
% Equivalent statement via elasticity. Let
% epsilon = -p D_p / D. Then
%   d epsilon / dA = -(p / D) [D_{pA} - D_p D_A / D].
% So d epsilon / dA < 0 iff D_{pA} > D_p D_A / D. Since
% D_p < 0 and D_A > 0, the RHS is negative — ads reduce
% elasticity for almost any positive (or even weakly
% negative) D_{pA}.
%
% The proposition's threshold (eq:sign_condition) is
% STRONGER than just "ads reduce elasticity"; it requires
% the magnitude of D_{pA} to outweigh the "smaller-margin"
% drag from incomplete pass-through.

\begin{corollary}[Separable-demand benchmark]
\label{cor:separable}
Suppose $D(p, A) = f(p) \cdot g(A)$ for some $f, g > 0$
with $f' < 0$ and $g'' < 0$ (multiplicative separability).
Then $D_{pA} = f'(p)\, g'(A)$, which has sign of $f'(p)
\cdot g'(A) = (-) \cdot (+) < 0$. The sign condition in
Proposition~\ref{prop:sign} fails. Hence $dA^*/dt < 0$:
the firm reduces advertising in response to the tax.
The sb\_incidence model's Assumption~3 (per-customer
demand independent of ads) is the limiting case
$g(A) = A$ — separability holds, and the same
conclusion follows.
\end{corollary}

\begin{corollary}[Persuasive-advertising case]
\label{cor:persuasive}
If advertising reduces demand elasticity sufficiently,
\eqref{eq:sign_condition} holds and $dA^*/dt > 0$. A
sufficient (not necessary) functional form: $D(p, A) =
\phi(A) \cdot D_0(p / \psi(A))$ with $\psi' > 0$, so
that higher $A$ shifts demand outward and rotates it
flatter in proportional space. Equivalently, ads act
as a brand-loyalty multiplier on the consumer's
reservation price.
\end{corollary}

% [SCAFFOLDING: empirical implication]
% Empirical implications to highlight:
%
% (i) The sb_incidence finding (sportsbook ads fall when
% GGR tax rises) is consistent with separable demand:
% sportsbook demand is dominated by extensive-margin
% participation, and ads acquire customers but do not
% meaningfully shift their per-bet wagering. Assumption 3
% in that paper is the right empirical call.
%
% (ii) The SSB result (Philly Spot TV ad spend rises
% relative to Pittsburgh by ~$11K/month after the
% beverage tax, see Section X) is consistent with
% Corollary~\ref{cor:persuasive}: in a CPG category where
% brands compete on persuasive recall and switching is
% one-aisle-over, ads plausibly rotate the demand curve.
% A positive D_{pA} is the canonical Becker-Murphy
% prediction.
%
% (iii) The same empirical exercise on liquor, beer,
% tobacco lets us test which side of the sign condition
% each industry sits on. Different signs across
% categories are themselves a finding.

\paragraph{Tax deductibility.}
% [SCAFFOLDING] Ads NOT tax-deductible — same setup as
% sb_incidence. With deductibility, the marginal cost
% of ads on the LHS of the ad FOC scales by (1-tau) for
% revenue taxes, but for a specific tax with non-
% deductible ads the algebra above is unchanged.

\paragraph{Bottom line.}
% [SCAFFOLDING] The sign of $dA^*/dt$ is the empirical
% test of whether ads in a given product category act as
% (a) attention-grabbers that don't affect price
% sensitivity (separable demand → ads fall with tax),
% or (b) persuasive demand-curve rotators (non-separable
% with D_{pA} > 0 → ads can rise with tax). The
% sb_incidence sportsbook result is the (a) case; the
% Philly SSB result is consistent with the (b) case.
% Stacking many tax events across categories lets us
% read off where each industry sits.
